\chapter{Chapter05 — Chow Theorem}
%%% generated by the blueprint pipeline from FABL.Chapter05.ChowTheorem
\section{Overview}

Declarations of the Lean module \texttt{FABL.Chapter05.ChowTheorem}.

\section{Declarations}

\begin{theorem}[Pointwise threshold representation inequality]
\lean{FABL.thresholdRepresentation_pointwise}
\label{FABL.thresholdRepresentation_pointwise}
\leanok
\uses{FABL.BooleanFunction, FABL.BooleanFunction.toReal, FABL.IsPolynomialThresholdRepresentation, FABL.thresholdSign_of_neg, FABL.thresholdSign_of_nonneg}
\difficulty{5}
Let $n$ be a natural number, let $f$ and $g$ be functions from $\\{-1,1\\}^n$ to
$\\{-1,1\\}$, and let $p:\\{-1,1\\}^n\\to\\mathbb R$ satisfy $f(x)=1$ when $p(x)\\ge 0$
and $f(x)=-1$ when $p(x)<0$ for every $x$. Suppose also that $p(x)\\ne 0$ for every $x$.
Then for every $x\\in\\{-1,1\\}^n$, $g(x)p(x)\\le f(x)p(x)$, and if $g(x)\\ne f(x)$,
then $g(x)p(x)<f(x)p(x)$.
\end{theorem}

\begin{theorem}[Eq of polynomial threshold representation of fourier coeff eq]
\lean{FABL.eq_of_polynomialThresholdRepresentation_of_fourierCoeff_eq}
\statementsource{Analysis-of-Boolean-Functions-by-Ryan-ODonnell}{pdf-ffbe9c7b3306-p116-b003}
\label{FABL.eq_of_polynomialThresholdRepresentation_of_fourierCoeff_eq}
\leanok
\uses{FABL.BooleanFunction, FABL.BooleanFunction.toReal, FABL.IsPolynomialThresholdRepresentation, FABL.fourierCoeff, FABL.fourierDegree, FABL.fourierDegree_le_iff, FABL.plancherel, FABL.thresholdRepresentation_pointwise, FABL.uniformInner}
\difficulty{5}
% dataset-pipeline: stub -- prose pending build_informal_statements.py
Statement pending regeneration.
\end{theorem}

\begin{theorem}[Polynomial threshold functions determined by low-degree coefficients]
\lean{FABL.eq_of_isPolynomialThreshold_of_fourierCoeff_eq}
\proofsource{Analysis-of-Boolean-Functions-by-Ryan-ODonnell}{pdf-ffbe9c7b3306-p116-b003}
\label{FABL.eq_of_isPolynomialThreshold_of_fourierCoeff_eq}
\leanok
\uses{FABL.BooleanFunction, FABL.BooleanFunction.toReal, FABL.IsPolynomialThreshold, FABL.eq_of_polynomialThresholdRepresentation_of_fourierCoeff_eq, FABL.exists_integer_polynomialThresholdRepresentation, FABL.fourierCoeff}
\difficulty{1}
Let $n$ and $k$ be nonnegative integers, and let $f,g:\\{-1,1\\}^n\\to\\{-1,1\\}$.
Suppose there is a real-valued function $p$ on the cube such that $f(x)$ is the sign of
$p(x)$ at every input, taking sign $1$ when $p(x)\\ge 0$, and the Fourier degree of $p$
is at most $k$. If the real-valued versions of $f$ and $g$ have equal Fourier
coefficients on every set of coordinates of size at most $k$, then $g=f$.
\end{theorem}

\begin{theorem}[Linear thresholds have degree at most one]
\lean{FABL.isPolynomialThreshold_one_of_isLinearThreshold}
\proofsource{Analysis-of-Boolean-Functions-by-Ryan-ODonnell}{
  pdf-ffbe9c7b3306-p113-b003,
  pdf-ffbe9c7b3306-p115-b005
}
\label{FABL.isPolynomialThreshold_one_of_isLinearThreshold}
\leanok
\uses{FABL.BooleanFunction, FABL.IsLinearThreshold, FABL.IsPolynomialThreshold, FABL.fourierDegree_affineSignLinearForm_le_one, FABL.signValue}
\difficulty{3}
Let $n$ be a nonnegative integer and let $f:\{-1,1\}^n\to\{-1,1\}$. If there are real
weights $a_1,\ldots,a_n$ and a real constant $a_0$ such that $f(x)$ is the sign of
$a_0+\sum_{i=1}^n a_i x_i$ for every $x$, with zero assigned the positive sign, then $f$
is represented by the sign of a real-valued polynomial on the cube whose Fourier degree
is at most $1$.
\end{theorem}

\begin{theorem}[Linear threshold functions determined by low-degree coefficients]
\lean{FABL.eq_of_isLinearThreshold_of_fourierCoeff_eq}
\proofsource{Analysis-of-Boolean-Functions-by-Ryan-ODonnell}{
  pdf-ffbe9c7b3306-p114-b002,
  pdf-ffbe9c7b3306-p114-b003
}
\label{FABL.eq_of_isLinearThreshold_of_fourierCoeff_eq}
\leanok
\uses{FABL.BooleanFunction, FABL.BooleanFunction.toReal, FABL.IsLinearThreshold, FABL.eq_of_isPolynomialThreshold_of_fourierCoeff_eq, FABL.fourierCoeff, FABL.isPolynomialThreshold_one_of_isLinearThreshold}
\difficulty{1}
Let $n$ be a nonnegative integer, and let $f,g:\\{-1,1\\}^n\\to\\{-1,1\\}$. Suppose $f$
is a linear threshold function: its value is the sign of an affine linear form in the
coordinates, with the convention that a nonnegative input has sign $1$. If the
real-valued versions of $f$ and $g$ have equal Fourier coefficients on every set of
coordinates of size at most one, then $g=f$.
\end{theorem}
